\documentclass[a4paper,twoside]{article}

\usepackage{morewrites}

\usepackage[svgnames,dvipsnames,x11names]{xcolor}
\usepackage{graphicx}
\usepackage[english]{babel}
\usepackage[colorlinks=true, linkcolor=NavyBlue, urlcolor=NavyBlue, citecolor=NavyBlue]{hyperref}
\usepackage[a4paper]{geometry}
\usepackage{ts-fonts}
\usepackage{csquotes}
\usepackage{amsmath}
\usepackage{float}
\usepackage{seqsplit}
\usepackage{ts-code}

\usepackage{lastpage}
\usepackage{refcount}
\makeatletter
\def\ps@plain{
  \let\@oddhead\@empty
  \let\@evenhead\@empty
  \def\@oddfoot{
    \hfil
    \ifnum\getpagerefnumber{LastPage}>1\relax
      \thepage
    \fi
    \hfil}
  \let\@evenfoot\@oddfoot
}
\makeatother

\title{Inline Math}
\author{}\date{}

\begin{document}

\maketitle

\thispagestyle{plain}
You can include inline math expressions using the standard \LaTeX{} delimiters \texttt{\textbackslash{}( ... \textbackslash{})} or \texttt{\$ ... \$} :

\begin{code}{markdown}{}{}
\begin{Verbatim}[commandchars=\\\{\},breaklines, breakanywhere, commandchars=\\\{\}]
The quadratic formula is given by \PYZbs{}(x = \PYZbs{}frac\PYZob{}\PYZhy{}b \PYZbs{}pm \PYZbs{}sqrt\PYZob{}b\PYZca{}2 \PYZhy{} 4ac\PYZcb{}\PYZcb{}\PYZob{}2a\PYZcb{}\PYZbs{})
or \PYZdl{}x = \PYZbs{}frac\PYZob{}\PYZhy{}b \PYZbs{}pm \PYZbs{}sqrt\PYZob{}b\PYZca{}2 \PYZhy{} 4ac\PYZcb{}\PYZcb{}\PYZob{}2a\PYZcb{}\PYZdl{}.
\end{Verbatim}
\end{code}
Rendered as:

\begin{displayquote}

The quadratic formula is given by $x = \frac{-b \pm \sqrt{b^2 - 4ac}}{2a}$
or $x = \frac{-b \pm \sqrt{b^2 - 4ac}}{2a}$.

\end{displayquote}
\section{Block Math}\label{block-math}

The Schrödinger equation in a non-relativistic case is written as:

$$
\imath \hbar \frac{\partial}{\partial t} \Psi(\mathbf{r},t) =
\left[ -\frac{\hbar^2}{2m} \nabla^2 + V(\mathbf{r},t) \right] \Psi(\mathbf{r},t)
$$

And the set of Maxwell's equations in differential form. The magnetic flux $\eqref{eq:max2}$ through any closed surface is zero, this implies that there are no magnetic monopoles.

\begin{align}
\nabla \cdot \vec{E} &= \frac{\rho}{\varepsilon_0} \quad &&\text{Gauss Law}\\[4pt]
\nabla \cdot \vec{B} &= 0 \quad &&\text{Gauss's law for electricity} \label{eq:max2}\\[4pt]
\nabla \times \vec{E} &= -\,\frac{\partial \vec{B}}{\partial t}
    \quad &&\text{Faraday's law}\\[4pt]
\nabla \times \vec{B} &= \mu_0 \vec{J} + \mu_0 \varepsilon_0
    \frac{\partial \vec{E}}{\partial t}
\quad &&\text{Ampère-Maxwell law}
\end{align}

\end{document}
